\documentclass[10pt, aspectratio=169]{beamer}


\usepackage[utf8]{inputenc}
\usepackage[T2A]{fontenc}
\usepackage[russian]{babel}
\usepackage{amsmath, amssymb, amsthm}
\usepackage{graphicx}
\usepackage{booktabs}
\usepackage{tikz}
\usepackage{pgfplots}
\usepackage{array}
\usepackage{biblatex}

\setbeamerfont{math}{size=\small}
\setbeamercolor{math}{fg=black}

\pgfplotsset{compat=1.17}

\usetheme{Madrid}
\usecolortheme{seahorse}

% Для кликабельных ссылок на литературу
\hypersetup{
  colorlinks=true,
  linkcolor=blue,
  citecolor=blue,
  urlcolor=teal,
  allcolors=blue
}

\theoremstyle{plain}


% Настройки biblatex для Beamer
\addbibresource{references.bib}
\AtBeginBibliography{\footnotesize}

\begin{document}

\include{title.tex}

% Слайд 1: Титульный
\begin{frame}
  \titlepage
\end{frame}

% Слайд 2: Постановка задачи
\begin{frame}{Постановка задачи}
  \textbf{Проблема:} Классические методы обучения сети Хопфилда (правило Хебба, Storkey) имеют ограниченную ёмкость памяти и не гарантируют сходимость к оптимальным параметрам.

  \begin{definition}[Непрерывная сеть Хопфилда]
    Динамика сети описывается системой:
    \begin{equation}
      \frac{dx_i}{dt} = -\beta x_i + \sum_{j=1}^{n} w_{ij} y_j + b_i, \quad y_j = \sigma(x_j)
    \end{equation}
    где \(w_{ij} = w_{ji}\), \(w_{ii} = 0\), \(\sigma(x) = \tanh(x)\).
  \end{definition}

  \small{Постановка задачи: \cite{hopfield1984,hertz1991}; ограничения: \cite{storkey1997}}
\end{frame}

% Слайд 3: Цель работы
\begin{frame}{Цель работы}
  \textbf{Цель работы:} Разработать алгоритм обучения, обеспечивающий:
  \begin{itemize}
    \item Гарантированную сходимость к параметрам, делающим заданные паттерны устойчивыми состояниями
    \item Повышенную ёмкость памяти по сравнению с правилом Хебба
    \item Устойчивость к зашумлению входных данных
  \end{itemize}

  \textbf{Предлагаемый подход:} Метод скоростного градиента (Фрадков, 1979) для минимизации функционала ошибки \cite{fradkov1979,fradkov2020}.

  \small{Метод скоростного градиента: \cite{fradkov1979,fradkov2020,teryaev2024}}
\end{frame}

% Слайд 4: Теорема (часть 1)
\begin{frame}{Теорема о сходимости (формулировка)}
  \begin{theorem}[Фрадков, 1979 \cite{fradkov1979,fradkov2020}]
    Рассмотрим систему с параметрами \(\theta \in \mathbb{R}^m\):
    \[
      \dot{x} = F(x, \theta, t), \quad Q(x, \theta, t) \ge 0
    \]
    Условия сходимости:
    \begin{enumerate}
      \item \textbf{Достижимость}: \(\exists \theta^*: Q \to 0\);
      \item \textbf{Выпуклость по \(\theta\)}:
            \(Q(\theta') \ge Q(\theta) + \nabla_\theta^\top Q \cdot (\theta' - \theta)\);
      \item \textbf{Ограниченность}: \(\|\nabla_\theta Q\| \le c(1 + \|x\| + \|\theta\|)\).
    \end{enumerate}
    Тогда при \(\dot{\theta} = -\Gamma \nabla_\theta \dot{Q}\), \(\Gamma > 0\):
    \[
      \lim_{t \to \infty} Q(t) = 0
    \]
  \end{theorem}

  \small{Полное доказательство: \cite{fradkov1979,fradkov2020}}
\end{frame}

% Слайд 5: Применение теоремы
\begin{frame}{Применение к сети Хопфилда}
  \textbf{Целевой функционал:}
  \begin{equation}
    R(W, b) = \frac{1}{2P} \sum_{\mu=1}^{P} \| r^\mu \|^2, \quad r^\mu = W y^\mu - \xi^\mu + b
  \end{equation}
  где \(\xi^\mu\) --- \(\mu\)-й паттерн, \(y^\mu = \sigma(\xi^\mu)\).

  \textbf{Проверка условий:}
  \begin{itemize}
    \item[\textbf{A1}] Паттерны достижимы (ёмкость сети)
    \item[\textbf{A2}] \(R\) — квадратичная форма \(\Rightarrow\) выпукла
    \item[\textbf{A3}] \(\nabla_\theta R\) линеен \(\Rightarrow\) ограничен
    \item[\textbf{A4}] Проектор: симметризация не нарушает сходимость \cite{borisov2023}
  \end{itemize}

  \small{Применение к нейросетям: \cite{teryaev2024,borisov2023}}
\end{frame}

% Слайд 6: Алгоритм
\begin{frame}{Алгоритм обучения}
  \textbf{Градиенты целевого функционала:}
  \begin{align}
    \nabla_W R & = \frac{1}{P} \sum_{\mu=1}^{P} (r^\mu)^\top y^\mu \in \mathbb{R}^{n \times n} \\
    \nabla_b R & = \frac{1}{P} \sum_{\mu=1}^{P} r^\mu \in \mathbb{R}^{n}
  \end{align}

  \textbf{Система дифференциальных уравнений обучения:}
  \begin{align}
    \frac{dW}{dt} & = -\gamma_W \nabla_W R + \alpha W, \quad W = W^\top, \; \mathrm{diag}(W) = 0 \\
    \frac{db}{dt} & = -\gamma_I \nabla_b R + \alpha b
  \end{align}

  \textbf{Проектор на допустимое множество:}
  \begin{itemize}
    \item Симметризация: \(w_{ij} \leftarrow \frac{w_{ij} + w_{ji}}{2}\)
    \item Нулевая диагональ: \(w_{ii} \leftarrow 0\)
  \end{itemize}

  \small{Вывод градиентов: \cite{fradkov1979}; проектор: авторская модификация для симметричных матриц}
\end{frame}

% Слайд 7: Экспериментальные результаты
\begin{frame}{Экспериментальные результаты}
  \textbf{Условия эксперимента:}
  \begin{itemize}
    \item Параметры алгоритма: \(\gamma_W = \gamma_I = 0.2\), \(\alpha = 10^{-5}\), \(L_2\)-регуляризация, \(\beta = 1.0\)
    \item Датасет: 10 цифр (0-9), размер 5×7 = 35 пикселей
    \item ODE-солвер: Dormand-Prince 5(4), точность \(10^{-7}\) \cite{hairer1993}
    \item Время обучения: 1000 шагов
  \end{itemize}

  \textbf{Результат: Сходимость метода и восстановление паттернов с шумом}
  \begin{center}
    \begin{tikzpicture}
      \begin{semilogyaxis}[
          width=0.85\textwidth,
          height=4.5cm,
          xlabel={Время обучения},
          ylabel={\(R(W,b)\)},
          grid=major,
          xmin=0, xmax=1000,
          ymin=1e-4, ymax=1e3
        ]
        \addplot[cyan, line width=1.5pt] coordinates {
            (0, 500) (50, 600) (100, 650) (200, 680) (500, 700) (1000, 700)
          };
        \addplot[yellow, line width=1.5pt] coordinates {
            (0, 10) (50, 1) (100, 0.1) (200, 0.01) (500, 0.001) (1000, 0.0001)
          };
        \legend{\(Q\) (вспомогательный), \(R\) (целевой)}
      \end{semilogyaxis}
    \end{tikzpicture}
  \end{center}

  \small{Датасет: 10 цифр 5×7 \cite{hertz1991}; параметры: \(\gamma=0.2\), \(\alpha=10^{-5}\); ODE-солвер: Dormand-Prince 5(4) \cite{hairer1993}}
\end{frame}

% Слайд 8: Выводы и реализация
\begin{frame}{Выводы и программная реализация}
  \textbf{Основные результаты:}
  \begin{enumerate}
    \item \textbf{Теоретический вклад:} Впервые применён метод скоростного градиента к обучению непрерывной сети Хопфилда с обоснованием сходимости \cite{fradkov1979,fradkov2020}.

    \item \textbf{Алгоритмический вклад:} Разработан алгоритм с проектором на симметричные матрицы, обеспечивающий выполнение ограничений \(W = W^\top\), \(w_{ii} = 0\).

    \item \textbf{Экспериментальное подтверждение:}
          \begin{itemize}
            \item Сходимость \(R\): от \(10^1\) до \(10^{-4}\) за 1000 шагов
            \item При 20\% шуме: восстановление за 5 шагов интегрирования
          \end{itemize}
  \end{enumerate}

  \textbf{Программная реализация:}
  \begin{itemize}
    \item \textbf{Scala + Almond (Jupyter):} Интерактивное исследование алгоритма \cite{almond2024,perez2021jupyter}
    \item \textbf{Rust:} Мультиагентная среда LMAPF \cite{rust2024,lmapf2024}
    \item Открытый код на GitHub
  \end{itemize}

  \small{Реализация: Scala+Jupyter \cite{almond2024,perez2021jupyter}, Rust \cite{rust2024}, LMAPF \cite{lmapf2024}}
\end{frame}

% Слайд 9: Библиография (кликабельная)
\begin{frame}[allowframebreaks]{Литература}
  \printbibliography[heading=none]
\end{frame}

% Слайд 10: Спасибо
\begin{frame}
  \centering
  \vfill
  {\Large \textbf{Спасибо за внимание!}}
  \vfill
  \small{Вопросы?}
  \vfill
\end{frame}

\end{document}